\documentclass{article}
\usepackage{graphicx} 
\usepackage[utf8]{inputenc}
\usepackage[T2A]{fontenc}
\usepackage[serbian]{babel}
\usepackage[unicode]{hyperref}
\usepackage{tabularx}
\usepackage{xurl}

\hypersetup{colorlinks,citecolor=green,filecolor=green,linkcolor=blue,urlcolor=blue}

\begin{document}
\begin{titlepage}

\newcommand{\HRule}{\rule{\linewidth}{0.4mm}}

\center 
	
\textsc{\LARGE Математички факултет}\\[3cm] 
	
\textsc{\Large Семинарски рад}\\[0.1cm]
	
\textsc{\large из Увода у информатику}\\[0.5cm] 
	
\HRule\\[0.4cm]
	
{\LARGE\bfseries Муров закон и експоненцијални развој технологије}\\[0.2cm] 
	
\HRule\\[2cm]

\vspace{17\baselineskip}
		

	
	\begin{minipage}{0.4\textwidth}
		\begin{flushleft}
			\large
			\textit{Студент}\\
			Владимир Хелета 200/2025
		\end{flushleft}
	\end{minipage}
\hspace*{1cm}
	\begin{minipage}{0.4\textwidth}
		\begin{flushright}
			\large
			\textit{Професор}\\
			проф. Данијела Симић
		\end{flushright}
	\end{minipage}
	
\vfill\vfill\vfill\vfill 
	
{\large Београд, \today} 
	
\vfill 
	
\end{titlepage}

\tableofcontents
\newpage

\section{Увод}
Развој рачунарске технологије у последњих неколико деценија представља један од 
најбржих технолошких напретака у историји. Једна од најпознатијих опсервација која 
описује овај напредак је \textbf{Муров закон}.

Овај закон је први пут формулисао \textit{Гордон Мур}, суоснивач компаније Intel, 
1965. године. Он је приметио да се број транзистора на интегрисаним колима 
\textbf{удвостручује приближно сваке две године}.
Ова појава је омогућила драматичан раст процесорске снаге и пад цена рачунарских система.

\begin{figure}[h!]
\centering
\includegraphics[width=0.8\textwidth]{murov_zakon}
\label{fig:murov_zakon}
\end{figure}

Модерни рачунари, паметни телефони и сервери дугују велики део свог развоја 
управо овом експоненцијалном тренду у повећању броја транзистора.

\section{Основни појмови}

\subsection{Дефиниција Муровог закона}


Муров закон представља емпиријску опсервацију да се број транзистора на интегрисаном 
колу удвостручује у приближно једнаком временском интервалу, најчешће око две године.


Овај раст се може приближно описати експоненцијалном функцијом.

\subsection{Математички модел}

Број транзистора током времена може се описати формулом:

\[
N(t) = N_0 \cdot 2^{t/2}
\]

где је:

\begin{itemize}
\item $N(t)$ број транзистора у тренутку $t$
\item $N_0$ почетни број транзистора
\item $t$ време у годинама
\end{itemize}

Ова формула показује да се раст одвија \textit{експоненцијално}, што је знатно брже 
од линеарног раста.

\section{Историјски развој процесора}

Развој микропроцесора кроз време може се јасно видети у следећој табели.


\begin{table}[!ht]
\centering
\caption{Раст броја транзистора кроз време}
\label{tab:rast}
\renewcommand{\arraystretch}{0.9}
\begin{tabularx}{\textwidth}{|X|c|c|}
\hline
\textbf{Година} & \textbf{Процесор} & \textbf{Број транзистора} \\
\hline
1971   & Intel 4004 & 2300 \\
1993   & Pentium & 3.1 милиона \\
2010   & Intel i7 & 1.17 милијарди \\
2020  & Apple M1 & 16 милијарди \\
\hline
\end{tabularx}
\end{table}

\subsection{Кључни фактори развоја}

Развој технологије омогућен је комбинацијом више фактора:

\begin{enumerate}
\item минијатуризација транзистора
\item напредак у производним процесима
\item нови материјали у производњи чипова
\end{enumerate}

Поред нумерисаних листа, често се користе и описне листе.

\begin{description}
\item[Минијатуризација] смањивање физичке величине транзистора
\item[Интеграција] постављање већег броја компоненти на један чип
\item[Ефикасност] мања потрошња енергије
\end{description}

\section{Математичка својства}


Ако се број транзистора удвостручује у једнаком временском интервалу, онда раст броја 
транзистора следи експоненцијалну функцију.


\textit{Скица доказа:}  
У сваком кораку вредност се множи са константом 2, што је основна карактеристика 
експоненцијалног раста.
Ако се тренд описан Муровим законом настави довољно дуго, рачунарска снага система 
расте експоненцијално у односу на време.


Ово значи да се за релативно кратко време могу појавити огромни скокови у перформансама 
рачунарских система.

\section{Ограничења Муровог закона}

Иако је Муров закон дуго важио, данас постоје одређена ограничења:

\begin{itemize}
\item физичка ограничења величине транзистора
\item повећање температуре чипова
\item раст цене производње
\end{itemize}

Због тога се модерни развој све више ослања на:

\begin{itemize}
\item паралелно процесирање
\item графичке процесоре
\item специјализоване AI чипове
\end{itemize}

\section{Закључак}

Муров закон је током више деценија био један од најважнијих принципа који су 
описивали развој рачунарске технологије. Иако се данас све више приближавамо 
физичким ограничењима минијатуризације, идеја експоненцијалног технолошког 
напретка и даље остаје кључна за разумевање будућег развоја рачунара.

\textbf{Развој нових архитектура, квантних рачунара и вештачке интелигенције 
вероватно ће представљати следећи корак у еволуцији рачунарских система.}

\end{document}

\end{document}
